\documentclass{article}
\usepackage[a4paper, total={6.5in, 10in}]{geometry}
\setlength{\parindent}{0pt}
\usepackage{tcolorbox}
\tcbset{colback=white!94!black, colframe=black!60!white, coltitle = white, boxrule=0.8pt, arc=2mm}
\newtcolorbox{boxs}[2][]{title= #2,#1}
\usepackage{datetime2}
\usepackage[british]{babel}
\usepackage{amsfonts}
\usepackage{amsmath}

\title{\textbf{Anlaysis on the Real Line HW 1}}
\author{Robert O'Connell}
\date{\today}
\begin{document}
\maketitle
\vspace{5mm}

\subsection*{BS - Page 11 Exercise 5}
(a)
\\
\[
A_1 = \{2k \ | \ k \in \mathbb{N}\}
\]
\[
A_2 = \{3k \ | \ k \in \mathbb{N}\}
\]

\[
A_1 \cap A_2
= \{ n \in \mathbb{N} \mid n = 2k \text{ and } n = 3m \text{ for some } k,m \in \mathbb{N} \}.
\]
These are all of the numbers which are multiples of 2 and of 3, i.e. multiples of 6 \\
Hence
\[
A_1 \cap A_2 = \{6k \ | \ k \in \mathbb{N}\} = A_5
\]
\\

(b)
\\
\[
\cup \{A_n \ | \ n \in \mathbb{N}\} 
\]
Consider \(A_k\), the multiples of \(k+1\), note \(k+1 \in A_k\) \\
Then \(\forall \ k \in \mathbb{N}, k + 1 \in A_k\) \\
i.e. for all \(A_k, \ k \in \mathbb{N}\) there is a set which includes \(k+1\). \\
No set inlcudes 1, since 1 is not a multiple
of \(k+1 \ \forall \ k \in \mathbb{N}\)
\[\Rightarrow \cup \{A_n \ | \ n \in \mathbb{N}\} = \mathbb{N} \backslash {1}\]
\\

\[
\cap \{A_n \ | \ n \in \mathbb{N}\}
\]
Consider \(A_k\), the multiples of \(k+1\), note \(k \notin A_k\) \\
Then \(\forall \ k \in \mathbb{N}, k \notin A_k\) \\
i.e. for all \(A_k, \ k \in \mathbb{N}\) there is a set which does not include \(k\), so it can't be in the intersection
\[\Rightarrow \cap \{A_n \ | \ n \in \mathbb{N}\} = \emptyset\]



\subsection*{BS - Page 11 Exercise 12}
Let \(f\) be a function, \(f: A \rightarrow B\) and let \(E,F\) be subsets of \(A\), 
to show \(f(E \cup F) = f(E) \cup f(F)\), we must show that \(F(E \cup F) \subseteq f(E) \cup f(F) \) and 
\(F(E) \cup f(F) \subseteq f(E \cup F) \)
\\

Firstly, \\ 
Let \(y \in f(E \cup F)\), then \(\exists \ x \in E \cup F\) such that \(y = f(x)\) \\
Since \(x \in E \cup F\), \(x \in E\) or \(x \in F\) 
\begin{itemize}
    \item If \(x \in E \Rightarrow y \in f(E)\)
    \item If \(x \in F \Rightarrow y \in f(F)\)
\end{itemize}

Thus we see if \(x \in E \cup F \Rightarrow y \in f(E)\) or \(y \in f(F)\) \\
Then \(y \in f(E) \cup f(F)\) \\
Therefore, \(f(E \cup F) \subseteq f(E) \cup f(F)\)
\\

For the second part,
Let \(y \in f(E) \cup f(F)\), then
\begin{itemize}
    \item If \(y \in f(E) \Rightarrow \exists \ x \in E\) such that \(y = f(x)\)
    \item If \(y \in f(F) \Rightarrow \exists \ x \in F\) such that \(y = f(x)\)
\end{itemize}

Thus we see that if \(y \in f(E) \cup f(F) \Rightarrow \exists\) \(x \in E\) or \(x \in F\) such that \(y = f(x)\) \\
Then \(\exists \ x \in E \cup F\) such that \(y = f(x)\) \\
\(\Rightarrow y \in f(E \cup F)\) \\
Therefore \(f(E) \cup f(F) \subseteq f(E \cup F) \)
\\

Hence \(f(E \cup F) = f(E) \cup f(F)\)
\\

\vspace{7mm}

Now we show that \(f(E \cap F) \subseteq f(E) \cap f(F)\),
\\ 

Let \(y \in f(E\cap F) \Rightarrow \ \exists \ x \in E \cap F\)  such that \( f(x) = y\). 
\\
Thus \(x \in E \) and \( x \in F \Rightarrow y \in f(E)\) and \(y \in f(F)\)

\(\Rightarrow y \in f(E)\cap f(F)\)

Hence \(f(E \cap F) \subseteq f(E) \cap f(F)\)













\subsection*{BS - Page 16 Exercise 16}
To find all natural number such that \(n^2 < 2^n\)
\\

Let \(P(n)\) be the proposition \(n^2 < 2^n\) with \(n \in \mathbb{N}\), \\
A quick check of the first few propositions shows that \(P(1), P(5), P(6), P(7)\) are
true while \(P(2), P(3), P(4)\) are false. \\

Let us check if \(P(n)\) holds for \(n > 4\),

\begin{itemize}
    \item \(P(5)\) is true since, \(25 < 32\)
    \item Assume \(P(k)\) holds.
Check if \(P(k+1)\) holds with the assumption
\begin{align*}
    k^2 &< 2^{k} \\
    k^2 + 2k + 1 &< 2^k + 2k + 1 \\
    (k+1)^2 &< 2^k + 2k + 1 \\
\end{align*}
We observe that \(\forall \ k > 4, \quad 2k + 1 < 2^k\)* \\
Which we will also prove by induction, \\

Base case: \(7 < 2^7 = 128\) \\
Induction step: Assume this holds for \(k = m\)

i.e. \(2m+1<2^m\) \\
\begin{align*}
    2m+1 + 2 &< 2^m +2 \\
    2(m+1) &< 2^m + 2 < 2(2^m) \ \ \text{Since } x+2 < 2x \ \forall \ x> 4, x \in \mathbb{N} \\
    &< 2^{m+1}
\end{align*}
Hence,
\begin{align*}
    (k+1)^2 &< 2^k + 2^k \\
    (k+1)^2 &< 2(2^k) \\
    (k+1)^2 &< 2^{k+1}
\end{align*}

Thus by induction, we see that \(n^2 < 2^n\) for all \(n > 4, n \in \mathbb{N}\)
\end{itemize}







\subsection*{BS - Page 21 Exercise 11}
Q. Use induction to prove that if a set \(S\) has \(n\) elements then \(\mathcal{P}(S)\) has \(2^n\) elements.

\begin{itemize}
    \item Base Case: A set with no elements (the empty set)
\[
\mathcal{P}(\emptyset) = \{\emptyset\}
\]
\[
2^0 = 1
\]

    \item Induction Step: Assume our proposition holds for any set \(A\) with \(k\) elements, i.e. it's powerset contains \(2^k\) sets.
Now consider the set \( S = A \cup \{x\}\) which has \(k+1\) elements. 

There are two categories of the subsets of \(S\), those which contain \(x\) and
those which don't. Of the ones the don't, by assumption there are exactly \(2^k\) of these, since these are all the possible subsets of \(A\). 

The other subsets are of the form \(B \cup \{x\}\) where 
\(B \subset A\), we know that there is exactly \(2^k\) subsets of A, thus there are exactly \(2^k\) subsets of \(S\) of this form. 

\(|\mathcal{P}(S)| = 2^k + 2^k = 2^{k+1}\)

Thus, by induction we have proved that  if a set \(S\) has \(n\) elements then \(\mathcal{P}(S)\) has \(2^n\) elements.

\end{itemize}
\end{document}